\section{Minimal Cocycle Support}

Let
\[
S = \{ e \in E(\Gfifteen) : \sigfun(e) = -1 \}.
\]

The set $S$ records the negative edges of the signing defining the
signed $2$-lift $\Gthirty \to \Gfifteen$.

\subsection{Gauge reduction}

Two signings related by vertex switching define the same cohomology
class. Explicitly, for a function $\tau : V(\Gfifteen) \to \{\pm 1\}$,
the switched signing is given by
\[
\sigfun'(uv) = \tau(u)\,\sigfun(uv)\,\tau(v).
\]

Thus the support $S$ is defined only up to switching equivalence.

\subsection{Minimal support}

We exhaustively search the switching orbit of $\sigfun$, fixing one
vertex to remove the global symmetry.

The computation shows that:
\begin{itemize}
\item the support can be reduced to size $6$,
\item no representative with fewer than $6$ negative edges exists,
\item there are $40$ distinct supports of size $6$ within the
gauge-fixed search.
\end{itemize}

This establishes that the minimal weight of the cohomology class is $6$.

\subsection{Explicit representative}

The table below records one representative of minimal support.

\begin{tabular}{c}
Edges in minimal support \\
\hline
$(0,1)$ \\
$(0,8)$ \\
$(2,13)$ \\
$(3,13)$ \\
$(10,11)$ \\
$(10,12)$ \\
\end{tabular}


\subsection{Interpretation}

Since the minimal support is nonzero, the cocycle defined by $\sigfun$
is nontrivial in $H^1(\Gfifteen; \mathbb{Z}_2)$.

The existence of multiple minimal representatives reflects the action
of the switching group on the cohomology class.
